\documentclass{article}

\usepackage[a4paper, left=1.5cm, right=1.5cm, top=2cm, bottom=2cm]{geometry}

\usepackage{../../../../components/mathtex}
\usepackage{pgfplots}
\usepackage{hyperref}
\usepackage{fancyhdr}
\usepackage{ccicons}
\usepackage{caption}
\usepackage{tabularx}
\usepackage{subcaption}

% Configuration des en-têtes et pieds de page
\pagestyle{fancy}
\fancyhf{} % reset tout

\fancyhead[L]{DL3 Math-Info PR5}
\fancyhead[C]{Intégration et probabilités}
\fancyhead[R]{2026-2027}

\fancyfoot[L]{Ewen Rodrigues de Oliveira\\\ccbyncsa}
\fancyfoot[R]{\thepage}

\begin{document}

\docTitle{Chapitre 3 : Fonctions de répartition et variables aléatoires}

\section{Fonction de répartition}

\subsection{Généralités et premières propriétés}
\definition{
    Soit $\mathbb{P}$ une mesure de probabilité sur $(\mathbb{R}, \mathcal{B}(\mathbb{R}))$. 
    On définit la \textbf{fonction de répartition} $F_\mathbb{P} : \mathbb{R} \to [0,1]$ par : 
    $$F_\mathbb{P}(x) = \mathbb{P}(]-\infty, x])$$
}

\theorem{Théorème}{Caractérisation de l'égalité des mesures}{false}{
    Soient $\mathbb{P}$ et $\mathbb{Q}$ deux mesures de probabilité sur $(\mathbb{R}, \mathcal{B}(\mathbb{R}))$. Alors :
    $$\mathbb{P} = \mathbb{Q} \iff F_\mathbb{P} = F_\mathbb{Q}$$
}
\noindent{\textbf{Preuve :}
    \begin{itemize}
        \item[$\Rightarrow$/] Direct.
        \item[$\Leftarrow$] $\{]-\infty, x], x \in \mathbb{R}\}$ est une classe de parties de $\mathbb{R}$ stable par intersection finies, qui engendre la tribu des boréliens, donc on peut utiliser le lemme d'unicité des mesures.
    \end{itemize}
}


\theorem{Proposition}{Croissance}{false}{
    Soit $\mathbb{P}$ une mesure de probabilité sur $(\mathbb{R}, \mathcal{B}(\mathbb{R}))$. Alors $F_\mathbb{P}$ est croissante.
}
\noindent{\textbf{Preuve :} Si $a < b$, $]-\infty, a] \subset ]-\infty, b]$ et donc
    \[
    F_{\mathbb{P}}(a) = \mathbb{P}(]-\infty, a]) \leq \mathbb{P}(]-\infty, b]) = F_{\mathbb{P}}(b)
    \]

où, pour la dernière égalité on a utilisé la propriété de limite décroissante pour une mesure de probabilités.}

\theorem{Proposition}{Continuité à droite}{false}{
    Soit $\mathbb{P}$ une mesure de probabilité sur $(\mathbb{R}, \mathcal{B}(\mathbb{R}))$. Alors $F_\mathbb{P}$ est continue à droite.
}
\noindent{\textbf{Preuve :} On a $]-\infty, a] = \bigcap_{n \in \mathbb{N}^*} ]-\infty, a + \frac{1}{n}]$, et en utilisant que $F_{\mathbb{P}}$ est croissante pour la première égalité ci-dessous,
    \begin{align*}
    \lim_{x \to a, x > a} F_{\mathbb{P}}(x) &= \lim_{n \to \infty} F_{\mathbb{P}}\left(a + \frac{1}{n}\right) \\
    &= \lim_{n \to \infty} \mathbb{P}\left(]-\infty, a + \frac{1}{n}]\right) = \mathbb{P}(]-\infty, a])
    \end{align*}}


\theorem{Proposition}{Limites aux bornes}{false}{
    Soit $\mathbb{P}$ une mesure de probabilité sur $(\mathbb{R}, \mathcal{B}(\mathbb{R}))$. Alors :
    $$\lim_{x \to -\infty} F_\mathbb{P}(x) = 0$$
    $$\lim_{x \to +\infty} F_\mathbb{P}(x) = 1$$
}

\theorem{Proposition}{Fonction de répartition sur un singleton}{false}{
    Soit $\mathbb{P}$ une mesure de probabilité sur $(\mathbb{R}, \mathcal{B}(\mathbb{R}))$. Il vient 
    $$\mathbb{P}(\{a\}) = F_\mathbb{P}(a) - \lim_{x \to a^-} F_\mathbb{P}(x)$$
}

\remark{En particulier, $F_\mathbb{P}$ est continue en $a \in \mathbb{R}$ si et seulement si $\mathbb{P}(\{a\}) = 0$.}
\remark{$F_\mathbb{P}$ admet un nombre au plus dénombrable de discontinuités.}

\training{Démontrer les deux dernières propositions (limites aux bornes et fonction de répartition sur un singleton).}
\noindent\carreaux{10}

\example{Soit $\mathbb{P}$ telle que \[
F_{\mathbb{P}}(x) = \sum_{k \in \mathbb{N}} \frac{2^k \cdot \exp(-3)}{k!} \indic{\{x \ge k\}} + (1 - \exp(-1)) * (1 - \exp(-x)) \indic{\{x \ge 0\}}
\]

\begin{itemize}
    \item \textbf{Que vaut $\mathbb{P}(\{0\})$ ?} On a :
    $\mathbb{P}(\{0\}) = F_\mathbb{P}(0) - \lim_{x \to 0^-} F_\mathbb{P}(x) = F_\mathbb{P}(0) - 0 = \frac{2^0 \cdot \exp(-3)}{0!} + (1 - \exp(-1)) * (1 - \exp(-0)) = \exp(-3)$
    \item \textbf{Déterminer $\{x \in \mathbb{R} : \mathbb{P}(\{x\}) > 0\}$.} Les seuls points où $\mathbb{P}(\{x\}) > 0$ sont les points $x \in \mathbb{N}$, car pour ces points, $\mathbb{P}(\{x\}) = \frac{2^x \cdot \exp(-3)}{x!}$. Ainsi, $\{x \in \mathbb{R} : \mathbb{P}(\{x\}) > 0\} = \mathbb{N}$.
    \item \textbf{Que vaut $\mathbb{P}([1, 2])$ ?} 

La mesure d'un intervalle fermé $[a, b]$ s'exprime par :
\[
\mathbb{P}([1, 2]) = F_\mathbb{P}(2) - \lim_{x \to 1^-} F_\mathbb{P}(x)
\]

Calculons séparément chaque terme :
\begin{itemize}
    \item Pour $F_\mathbb{P}(2)$, les termes $k \in \{0, 1, 2\}$ contribuent à la somme :
    \[
    F_\mathbb{P}(2) = e^{-3}\left(\frac{2^0}{0!} + \frac{2^1}{1!} + \frac{2^2}{2!}\right) + (1 - e^{-1})(1 - e^{-2}) = 5e^{-3} + (1 - e^{-1})(1 - e^{-2})
    \]
    \item Pour la limite à gauche $\lim_{x \to 1^-} F_\mathbb{P}(x)$, seul le terme $k=0$ intervient (car $x < 1$) :
    \[
    \lim_{x \to 1^-} F_\mathbb{P}(x) = e^{-3}\left(\frac{2^0}{0!}\right) + (1 - e^{-1})(1 - e^{-1}) = e^{-3} + (1 - e^{-1})^2
    \]
\end{itemize}

En faisant la différence, on obtient :
\[
\mathbb{P}([1, 2]) = (5e^{-3} - e^{-3}) + (1 - e^{-1})\Big[(1 - e^{-2}) - (1 - e^{-1})\Big]
\]
\[
\mathbb{P}([1, 2]) = 4e^{-3} + (1 - e^{-1})(e^{-1} - e^{-2}) = 4e^{-3} + e^{-1}(1 - e^{-1})^2
\]
\end{itemize}
}

\remark{Cette fonction de répartition est un mélange d'une loi de Poisson (avec paramètre $\lambda = 3$) et d'une loi exponentielle (avec paramètre $\lambda = 1$).}

\theorem{Théorème}{}{false}{
    Soit $F : \mathbb{R} \to \mathbb{R}$ croissante, continue à droite, telle que $\lim_{x \to -\infty} F(x) = 0$ et $\lim_{x \to +\infty} F(x) = 1$. Alors il existe une unique mesure de probabilité $\mathbb{P}$ sur $(\mathbb{R}, \mathcal{B}(\mathbb{R}))$ telle que $F = F_\mathbb{P}$.
}

\section{Variables aléatoires}

\subsection{Définition}
\definition{
    Soit $(\Omega, \mathcal{F}, \mathbb{P})$ un espace de probabilité. Soit $E$ muni de la tribu $\mathcal{E}$.
    Toute application mesurable $X : \Omega \to E$ est une variable aléatoire à valeurs dans $E$.\\\\
    La loi de $X$, $\mathbb{P}_X$, est la mesure-image de $\mathbb{P}$ par $X$.
}

\subsection{Variable aléatoire réelle}
\subsubsection{Propriétés immédiates}
\remark{Si $E = \mathbb{R}$, on parle de variable aléatoire réelle (v.a.r.). Sa loi $\mathbb{P}_X$ est une mesure de probabilité sur $(\mathbb{R}, \mathcal{B}(\mathbb{R}))$.}

\theorem{Proposition}{Caractérisation par la fonction de répartition}{false}{
    Soit $X$ une variable aléatoire réelle. Alors sa loi $\mathbb{P}_X$ est entièrement caractérisée par sa fonction de répartition $F_X = F_{\mathbb{P}_X}$, définie par :
    $$F_X(a) = \mathbb{P}_X(]-\infty, a]) = \mathbb{P}(X \le a)$$
}

\theorem{Proposition}{Image par une fonction borélienne}{false}{
    Soit $X$ une variable aléatoire réelle et $f : \mathbb{R} \to \mathbb{R}$ une fonction borélienne. Alors $f(X)$ est une variable aléatoire réelle.
}

\theorem{Proposiiton}{Couple de variables aléatoires}{false}{
    $(X,Y)$ est une variable aléatoire à valeurs dans $\mathbb{R}^2$ si et seulement si $X$ et $Y$ sont des variables aléatoires réelles sur le même espace.\\
    Dans ce cas, $X+Y$ et $X \cdot Y$ sont des variables aléatoires réelles.
}

\theorem{Proposition}{Suite de variables aléatoires}{false}{
    Soit $(X_n)_{n \in \mathbb{N}}$ une suite de variables aléatoires réelles sur un même espace. Alors $\sup_{n \in \mathbb{N}} X_n$, $\inf_{n \in \mathbb{N}} X_n$, $\limsup_{n \to \infty} X_n$ et $\liminf_{n \to \infty} X_n$ sont des variables aléatoires réelles (sous réserve d'existence).
}

\subsubsection{Espérance et variance}

\definition{
    Soit $X$ est une v.a.r. positive ou dans $\mathcal{L}^1(\Omega, \mathcal{F}, \mathbb{P})$. On définit l'espérance de $X$ par :
    $$\mathbb{E}[X] = \int_\Omega X d\mathbb{P}$$
}

\definition{
    Soit $X$ est dans $\mathcal{L}^2(\Omega, \mathcal{F}, \mathbb{P})$. On définit la variance de $X$ par :
    $$\text{Var}(X) = \mathbb{E}[(X - \mathbb{E}[X])^2] = \begin{cases}
\mathbb{E}[X^2] - (\mathbb{E}[X])^2 & \text{si } \mathbb{E}[|X|] < \infty \\
\end{cases}$$
}

\theorem{Théorème de transfert}{}{false}{
    Soit $phi : \mathbb{R} \to \mathbb{R}$ une fonction borélienne et $X$ une variable aléatoire réelle. Alors :
    $$\mathbb{E}[\phi(X)] = \int_\Omega \phi(X(\omega)) d\mathbb{P}(\omega) = \int_\mathbb{R} \phi(x) d\mathbb{P}_X(x)$$
    Si $\phi : \mathbb{R} \to \mathbb{R}$ est borélienne, alors $\phi \circ X \in \mathcal{L}^1(\Omega, \mathcal{F}, \mathbb{P})$ si et seulement si $\phi \in \mathcal{L}^1(\mathbb{R}, \mathcal{B}(\mathbb{R}), \mathbb{P}_X)$ et dans ce cas :
    $$\mathbb{E}[\phi(X)] = \int_\Omega \phi(X(\omega)) d\mathbb{P}(\omega) = \int_\mathbb{R} \phi(x) d\mathbb{P}_X(x)$$
}

\subsubsection{Lois usuelles}

\definition{
    Soit $X$ une variable aléatoire réelle. On dit que $X$ suit la loi de Bernoulli de paramètre $p \in [0,1]$ si :
    $$\mathbb{P}(X = 1) = p \quad \text{et} \quad \mathbb{P}(X = 0) = 1 - p$$
    Autrement dit $X \sim Ber(p) \iff \mathbb{P}_X = p \delta_1 + (1 - p) \delta_0 \iff F_X(x) = (1-p) \indic{x \ge 0} + p \indic{x \ge 1}$
}
\remark{Pour $\phi : \mathbb{R} \to \mathbb{R}$ borélienne positive, on a $\mathbb{E}[\phi(X)] = (1-p) \phi(0) + p \phi(1)$. En particulier, $\mathbb{E}[X] = p$ et $\text{Var}(X) = p(1-p)$.}

\definition{
    Soit $n \in \mathbb{N}^*$, $p \in [0,1]$. On dit que $X$ suit la loi binomiale de paramètres $n$ et $p$ et on note $X \sim Bin(n,p)$ si :
    $$\mathbb{P}(X = k) = \binom{n}{k} p^k (1-p)^{n-k}, \quad k = 0, \ldots, n$$
    Autrement dit, $X \sim Bin(n,p) \iff \mathbb{P}_X = \sum_{k=0}^n \binom{n}{k} p^k (1-p)^{n-k} \delta_k \iff F_X(x) = \sum_{k=0}^{n} \binom{n}{k} p^k (1-p)^{n-k} \indic{\{x \ge k\}}$
}
\remark{Pour $\phi : \mathbb{R} \to \mathbb{R}$ borélienne positive, on a $\mathbb{E}[\phi(X)] = \sum_{k=0}^n \binom{n}{k} p^k (1-p)^{n-k} \phi(k)$. En particulier, $\mathbb{E}[X] = np$ et $\text{Var}(X) = np(1-p)$.}

\definition{
    Soit $p \in ]0,1]$. On dit que $X$ suit la loi géométrique de paramètre $p$ et on note $X \sim Geom(p)$ si :
    $$\mathbb{P}(X = k) = p (1-p)^{k-1}, \quad k \in \mathbb{N}^*$$
    Autrement dit, $X \sim Geom(p) \iff \mathbb{P}_X = \sum_{k=1}^{\infty} p (1-p)^{k-1} \delta_k \iff F_X(x) = \sum_{k=1}^{+\infty} p (1-p)^{k-1} \indic{\{x \ge k\}}$
}
\remark{Pour $\phi : \mathbb{R} \to \mathbb{R}$ borélienne positive, on a $\mathbb{E}[\phi(X)] = \sum_{k=1}^{\infty} p (1-p)^{k-1} \phi(k)$.}
\remark{En particulier, $\mathbb{E}[t^X] = \frac{p t}{1 - (1-p) t}$ et on en déduit $\mathbb{E}[X] = \frac{1}{p}$ et $\text{Var}(X) = \frac{1-p}{p^2}$.}

\definition{
    Soit $\lambda > 0$. On dit que $X$ suit la loi de Poisson de paramètre $\lambda$ et on note $X \sim Poi(\lambda)$ si :
    $$\mathbb{P}(X = k) = \frac{\lambda^k e^{-\lambda}}{k!}, \quad k \in \mathbb{N}$$
    Autrement dit, $X \sim Poi(\lambda) \iff \mathbb{P}_X = \sum_{k=0}^{\infty} \frac{\lambda^k e^{-\lambda}}{k!} \delta_k \iff F_X(x) = \sum_{k=0}^{+\infty} \frac{\lambda^k e^{-\lambda}}{k!} \indic{\{x \ge k\}}$
}
\remark{Pour $\phi : \mathbb{R} \to \mathbb{R}$ borélienne positive ou telle que $\mathbb{E}[|\phi(X)|] < \infty$, on a $\mathbb{E}[\phi(X)] = \sum_{k=0}^{\infty} \frac{\lambda^k e^{-\lambda}}{k!} \phi(k)$. En particulier, $\mathbb{E}[X] = \lambda$ et $\text{Var}(X) = \lambda$.}

\definition{
    Soit $n \in \mathbb{N}^*$. On dit que $X$ suit la loi uniforme sur $\{1, \ldots, n\}$ et on note $X \sim U(\{1, \ldots, n\})$ si :
    $$\mathbb{P}(X = k) = \frac{1}{n}, \quad k = 1, \ldots, n$$
    Autrement dit, $X \sim U(\{1, \ldots, n\}) \iff \mathbb{P}_X = \sum_{k=1}^{n} \frac{1}{n} \delta_k \iff F_X(x) = \sum_{k=1}^{n} \frac{1}{n} \indic{\{x \ge k\}}$
}
\remark{Pour $\phi : \mathbb{R} \to \mathbb{R}$ borélienne ou telle que $\mathbb{E}[|\phi(X)|] < \infty$, on a $\mathbb{E}[\phi(X)] = \sum_{k=1}^{n} \frac{1}{n} \phi(k)$. En particulier, $\mathbb{E}[X] = \frac{n+1}{2}$ et $\text{Var}(X) = \frac{n^2 - 1}{12}$.}

\definition{
    On dit que $X$ suit la loi uniforme sur $[a, b]$ et on note $X \sim U([a, b])$ si :
    $$F_X(t) = \begin{cases}
        0 & \text{si } t < a, \\
        \frac{t - a}{b - a} & \text{si } a \le t \le b, \\
        1 & \text{si } t > b.
    \end{cases}$$
    On dit aussi que $X$ (on devrait dire $\mathbb{P}_X$) possède la densité $f_X(t) = \frac{1}{b - a} \indic{[a, b]}(t)$ vis-à-vis de la mesure de Lebesgue.
}  
\remark{Pour $\phi : \mathbb{R} \to \mathbb{R}$ borélienne ou telle que $\mathbb{E}[|\phi(X)|] < \infty$, on a $\mathbb{E}[\phi(X)] = \int_a^b \frac{\phi(t)}{b - a} dt$. En particulier, $\mathbb{E}[X] = \frac{a + b}{2}$ et $\text{Var}(X) = \frac{(b - a)^2}{12}$.}
\ndlr{vérifier la remarque}

\definition{
    Soient $\lambda > 0$. On dit que $X$ suit la loi exponentielle de paramètre $\lambda$ et on note $X \sim Exp(\lambda)$ si :
    $$F_X(t) = (1 - e^{-\lambda t}) \indic{t \ge 0} = \int_{-\infty}^t \lambda e^{-\lambda x} \indic{x \ge 0} dx$$
    On dit que $X$ possède la densité $f_x = \lambda e^{-\lambda x} \indic{x \ge 0}$ vis-à-vis de la mesure de Lebesgue.
}
\remark{Pour $\phi : \mathbb{R} \to \mathbb{R}$ borélienne ou telle que $\mathbb{E}[|\phi(X)|] < \infty$, on a $\mathbb{E}[\phi(X)] = \int_0^{\infty} \lambda e^{-\lambda x} \phi(x) dx$. En particulier, $\mathbb{E}[X] = \frac{1}{\lambda}$ et $\text{Var}(X) = \frac{1}{\lambda^2}$.}

\definition{
    On dit que $X$ suit une loi normale centrée réduite et on note $X \sim N(0, 1)$ si :
    $$F_X(t) = \frac{1}{\sqrt{2\pi}} \int_{-\infty}^t e^{-\frac{x^2}{2}} dx$$
    On dit que $X$ possède la densité $f_X(x) = \frac{1}{\sqrt{2\pi}} e^{-\frac{x^2}{2}}$ vis-à-vis de la mesure de Lebesgue.
}
\remark{Pour $\phi : \mathbb{R} \to \mathbb{R}$ borélienne positive ou telle que $\mathbb{E}[|\phi(X)|] < \infty$, on a $\mathbb{E}[\phi(X)] = \int_{-\infty}^{\infty} \frac{\phi(x)}{\sqrt{2\pi}} e^{-\frac{x^2}{2}} dx$. En particulier, $\mathbb{E}[X] = 0$ et $\text{Var}(X) = 1$.}

\definition{
    Soient $\mu \in \mathbb{R}$ et $\sigma > 0$. On dit que $X$ suit une loi normale de paramètres $\mu$ et $\sigma^2$ et on note $X \sim N(\mu, \sigma^2)$ si :
    $$F_X(t) = \frac{1}{\sigma \sqrt{2\pi}} \int_{-\infty}^t e^{-\frac{(x - \mu)^2}{2\sigma^2}} dx$$
    On dit que $X$ possède la densité $f_X(x) = \frac{1}{\sigma \sqrt{2\pi}} e^{-\frac{(x - \mu)^2}{2\sigma^2}}$ vis-à-vis de la mesure de Lebesgue.
}
\remark{Pour $\phi : \mathbb{R} \to \mathbb{R}$ borélienne positive ou telle que $\mathbb{E}[|\phi(X)|] < \infty$, on a $\mathbb{E}[\phi(X)] = \int_{-\infty}^{\infty} \frac{\phi(x)}{\sigma \sqrt{2\pi}} e^{-\frac{(x - \mu)^2}{2\sigma^2}} dx$. En particulier, $\mathbb{E}[X] = \mu$ et $\text{Var}(X) = \sigma^2$.}
\courseEnd{Fin du cours du 22 octobre 2026}

\definition{
    Soit $a > 0, \lambda > 0$. On dit que $X$ suit une loi Gamma de paramètres $a$ et $\lambda$ et on note $X \sim \Gamma(a, \lambda)$ si :
    $$F_X(t) = \int_{-\infty}^t \frac{\lambda^a}{\Gamma(a)} x^{a-1} e^{-\lambda x} \indic{x \ge 0} dx$$
    où $\Gamma(a) = \int_0^{\infty} x^{a-1} e^{-x} dx$ est la fonction Gamma. On dit que $X$ possède la densité $f_X(x) = \frac{\lambda^a}{\Gamma(a)} x^{a-1} e^{-\lambda x} \indic{x \ge 0}$ vis-à-vis de la mesure de Lebesgue.
}
\remark{Pour $\phi : \mathbb{R} \to \mathbb{R}$ borélienne positive ou telle que $\mathbb{E}[|\phi(X)|] < \infty$, on a $\mathbb{E}[\phi(X)] = \int_0^{\infty} \frac{\lambda^a}{\Gamma(a)} x^{a-1} e^{-\lambda x} \phi(x) dx$. En particulier, $\mathbb{E}[X] = \frac{a}{\lambda}$ et $\text{Var}(X) = \frac{a}{\lambda^2}$.}

\subsubsection{Tableau récapitulatif des lois usuelles}

\begin{table}[h!]
\centering
\renewcommand{\arraystretch}{2.0}
\setlength{\tabcolsep}{5pt} % On réduit légèrement l'espacement latéral
\small
\begin{tabularx}{\linewidth}{|l|X|X|c|c|}
\hline
\textbf{Loi discrète} & \textbf{Support / Loi} $\mathbb{P}(X = k)$ & \textbf{Fonction de Répartition} $F_X(x)$ & $\mathbb{E}[X]$ & $\text{Var}(X)$ \\
\hline
Bernoulli \newline {\scriptsize $\mathcal{B}er(p)$} & $\mathbb{P}(X=1)=p, \, \mathbb{P}(X=0)=1-p$ \newline ($p \in [0,1]$) & $(1-p)\indic{x \ge 0} + p\indic{x \ge 1}$ & $p$ & $p(1-p)$ \\
\hline
Binomiale \newline {\scriptsize $\mathcal{B}in(n,p)$} & $\binom{n}{k} p^k (1-p)^{n-k}, \quad k \in \{0, \ldots, n\}$ & $\sum_{k=0}^{n} \binom{n}{k} p^k (1-p)^{n-k} \indic{x \ge k}$ & $np$ & $np(1-p)$ \\
\hline
Géométrique \newline {\scriptsize $\mathcal{G}eom(p)$} & $p(1-p)^{k-1}, \quad k \in \mathbb{N}^*$ ($p \in ]0,1]$) & $\sum_{k=1}^{+\infty} p (1-p)^{k-1} \indic{x \ge k}$ & $\dfrac{1}{p}$ & $\dfrac{1-p}{p^2}$ \\
\hline
Poisson \newline {\scriptsize $\mathcal{P}oi(\lambda)$} & $\dfrac{\lambda^k e^{-\lambda}}{k!}, \quad k \in \mathbb{N}$ ($\lambda > 0$) & $\sum_{k=0}^{+\infty} \dfrac{\lambda^k e^{-\lambda}}{k!} \indic{x \ge k}$ & $\lambda$ & $\lambda$ \\
\hline
Uniforme \newline {\scriptsize $\mathcal{U}(\{1, \ldots, n\})$} & $\dfrac{1}{n}, \quad k \in \{1, \ldots, n\}$ & $\sum_{k=1}^{n} \dfrac{1}{n} \indic{x \ge k}$ & $\dfrac{n+1}{2}$ & $\dfrac{n^2 - 1}{12}$ \\
\hline
\end{tabularx}
\end{table}

\vspace{0.5cm}

\begin{table}[h!]
\centering
\renewcommand{\arraystretch}{2.2}
\setlength{\tabcolsep}{5pt}
\small
\begin{tabularx}{\linewidth}{|l|X|X|c|c|}
\hline
\textbf{Loi à densité} & \textbf{Densité} $f_X(x)$ & \textbf{Fonction de Répartition} $F_X(t)$ & $\mathbb{E}[X]$ & $\text{Var}(X)$ \\
\hline
Uniforme \newline {\scriptsize $\mathcal{U}([a,b])$} & $\dfrac{1}{b-a} \indic{[a,b]}(x)$ & $\begin{cases} 0 & t < a \\ \dfrac{t-a}{b-a} & a \le t \le b \\ 1 & t > b \end{cases}$ & $\dfrac{a+b}{2}$ & $\dfrac{(b-a)^2}{12}$ \\
\hline
Exponentielle \newline {\scriptsize $\mathcal{E}xp(\lambda)$} & $\lambda e^{-\lambda x} \indic{x \ge 0} \quad (\lambda > 0)$ & $(1 - e^{-\lambda t}) \indic{t \ge 0}$ & $\dfrac{1}{\lambda}$ & $\dfrac{1}{\lambda^2}$ \\
\hline
Normale \newline {\scriptsize $\mathcal{N}(0,1)$} & $\dfrac{1}{\sqrt{2\pi}} e^{-\frac{x^2}{2}}$ & $\dfrac{1}{\sqrt{2\pi}} \int_{-\infty}^t e^{-\frac{x^2}{2}} dx$ & $0$ & $1$ \\
\hline
Normale \newline {\scriptsize $\mathcal{N}(\mu, \sigma^2)$} & $\dfrac{1}{\sigma \sqrt{2\pi}} e^{-\frac{(x-\mu)^2}{2\sigma^2}} \quad (\sigma > 0)$ & $\dfrac{1}{\sigma \sqrt{2\pi}} \int_{-\infty}^t e^{-\frac{(x-\mu)^2}{2\sigma^2}} dx$ & $\mu$ & $\sigma^2$ \\
\hline
Gamma \newline {\scriptsize $\Gamma(a, \lambda)$} & $\dfrac{\lambda^a}{\Gamma(a)} x^{a-1} e^{-\lambda x} \indic{x \ge 0} \quad (a,\lambda > 0)$ & $\int_{-\infty}^t \dfrac{\lambda^a}{\Gamma(a)} x^{a-1} e^{-\lambda x} \indic{x \ge 0} dx$ & $\dfrac{a}{\lambda}$ & $\dfrac{a}{\lambda^2}$ \\
\hline
\end{tabularx}
\end{table}

\newpage
\tableofcontents
\section*{Crédits}
Ce cours est inspiré du cours de l'Université Paris Cité réalisés par Mathieu Merle. Assurez-vous de citer correctement la source si vous utilisez ce cours dans vos travaux.\\\\
Ce cours est mis à disposition selon les termes de la Licence Creative Commons \ccbyncsa\ \textbf{Attribution - Pas d'Utilisation Commerciale - Partage dans les Mêmes Conditions 4.0 International}. 
Pour voir une copie de cette licence, visitez \url{http://creativecommons.org/licenses/by-nc-sa/4.0/}.

\end{document}